% !TeX root = ../document_maitre.tex
% document de glossaire.


\newglossaryentry{ina}{
	name={inaliéniable},
	description={En terme juridique, cela désigne ce qui ne peut être retiré, altéré ni même vendu ou donné. Il ne peut être modifié}
}

\newglossaryentry{inv}{
	name={inviolable},
	description={En terme juridique, désigne un caractère protégé qu'on ne peut enfreindre. Il ne peut être enfreint et s'il est associé à un objet, cela rend ce dernier également inviolable.}
}

\newglossaryentry{rec}{
	name={Récolement décennale},
	description={Mesure imposée par l'État français contraignant les musées de France à vérifier l'état de leurs collections tous les dix ans. C'est une opération de vérifications des pièces présentes, absentes, de leurs états, localisations, marquage et de leur conformité vis-à-vis de son inscription à l'inventaire.}
}

\newglossaryentry{vand}{
	name={vandalisme},
	description={Inventé par l'Abbé Grégoire en 1792\footcite{sprigathVandalismeRevolutionnaire179217941980a} à l'occasion de son étude sur l'impact révolutionnaire sur les biens culturels français, il a assez peu changé depuis. Le vandalisme se définit comme tout acte de destruction ou de dégradation volontaire d'\oe{uvres} d'art, de monuments ou de matériel}
}

\newglossaryentry{imgani}{
	name={image animée},
	plural={images animées},
	description={Ce terme désigne tout objet constitué par une suite d'images capturées par différents procédés et dont la diffusion à la chaîne crée du mouvement. Par image animée, on désigne bien souvent les documents audivisuels comme les films, les séries, les documentaires, vidéos que ce soit amateur ou professionnel}
}

\newglossaryentry{BNF}{
	name={BNF},
	description={Successeure de la Bibliothèque Royale, la Bibliothèque Nationale de France, ou BNF, est l'institution en charge de la collecte, l'archivage et la préservation de l'édition en France notamment à travers le dépôt légal dont elle est la seule responsable. Elle a également à charge une des plus grandes collections patrimoniale française}
}

\newglossaryentry{IFLA}{
	name={IFLA},
	description={La Fédération internationale des associations et institutions de bibliothèques est une organisation internationale représentant les bibliothèques à l'échelle mondiale. Cette organisation se charge de défendre les intérêts de ces institutions et de développer la gestion des collections de bibliothèques de manière commune.\footnote{ Pour plus d'informations voir \url{https://www.ifla.org/}}}
}

\newglossaryentry{CIDOC}{
	name={CIDOC},
	description={Le Comité international pour la documentation est un comité international représentant les musées à l'échelle mondiale. Il se spécialise dans la rédaction et l'application de normes pour la gestion et la description des collections des musées.\footnote{Pour plus d'informations voir \url{https://cidoc.mini.icom.museum/}}}
}

\newglossaryentry{IA}{
	name={intelligence artificielle},
	text={IA},
	first={intelligence artificielle},
	plural={intelligences artificielles},
	description={Ensemble vaste de technologies et procédés informatiques et d'ingénierie qui visent à permettre à l'ordinateur d'effectuer des tâches assimilées à l'intelligence humaines.\footcites{IntelligenceArtificielleQuoi}{IntelligenceArtificielleDefinition2020} Ses origines remontent aux travaux de John McCarthy en 1956 et qui voit ses premières expériences concluantes apparaître dans les années 1990}
}

\newglossaryentry{bgdt}{
	name={\textit{Big Data}},
	sort={big data},
	description={Phénomène d'augmentation quantitative importantes des données dans tous les domaines à l'échelle mondiale à partir de 2010. On utilise ce terme pour faire référence aux nouvelles techniques de traitement de la données mises en place afin de répondre aux défis posés par ce phénomène}
}

\newglossaryentry{OCR}{
	name={\textit{Optical Character Recognition}},
	first={\textit{Optical Character Recognition}},
	text={OCR},
	sort={optical character recognition},
	description={L'\textit{optical character recognition} créé par l'ingénieur Autrichien Gustave Tauschek, ou reconnaissance optique de caractère en français, désigne la capacité conférée à un système informatique de reconnaître des mots typographiés. Le premier système à pouvoir le faire est la machine Gismo inventée en 1951 par David Shepard.}
}

\newglossaryentry{HTR}{
	name={\textit{Handwritten Text Recognition}},
	first={\textit{Handwritten Text Recognition}},
	text={HTR},
	sort={handwritten text recognition},
	description={L'\textit{Handwritten Text Recognition} est basé sur la même idée que l'OCR, la reconnaissance d'écriture manuscrite, ou HTR, applique le même principe à l'analyse des écritures manuscrites. On l'applique surtout aux documents anciens nécessitant des connaissances paléographiques.}
}

\newglossaryentry{NLP}{
	name={\textit{Natural Language Processing}},
	first={\textit{Natural Language Processing}},
	text={NLP},
	sort={natural language processing},
	description={Le \textit{Natural Language Processing}, ou traitement automatique du langage naturel, est un sous-domaine multidisciplinaire de l'intelligence artificielle, faisant notamment appel à l'ingénierie informatique et la linguistique, basé sur l'utilisation du \gls{ML} pour conférer à un système informatique la capacité d'analyser, comprendre et communiquer dans ce que l'on appelle le langage naturel ou langage courant.}
}

\newglossaryentry{CV}{
	name={\textit{computer vision}},
	sort={computer vision},
	description={La \textit{computer vision}, ou vision par ordinateur en français, est sous-domaine multidisciplinaire de l'\gls{IA}, cela désigne la capacité d'analyse et d'interprétation des images données à un ordinateur ou système informatique.}
}

\newglossaryentry{ML}{
	name={\textit{machine learning}},
	sort={machine learning},
	description={Domaine d'étude de l'intelligence artificielle qui vise à permettre à un système informatique basé, généralement, sur des réseaux de neurones artificiels, d'apprendre. Ce domaine se consacre à développer de nouvelles manières de conférer à l'intelligence artificielle des moyens d'apprendre de différentes manières.} 
}

\newglossaryentry{UNESCO}{
	name={UNESCO},
	description={L'Organisation des Nations unies pour l'éducation, la science et la culture est une organisation internationale créée en 1945 qui a pour but de promulguer la coopération internationale autour des champs de l'éducation, de la culture, de la science et de la communication. Ces missions sont plurielles, mais elle est très largement connue pour ses initiatives sur le patrimoine : reconnaissance du patrimoine immatériel en 2003 et la construction d'un patrimoine mondial commencée en 1972.} 
}

\newglossaryentry{DPL}{
	name={\textit{Deep Learning}},
	sort={deep learning},
	description={Branche du \gls{ML}, il désigne à la fois un modèle d'architecture de réseaux neuronaux d'\gls{IA} et une méthode d'entraîner les technologies basée sur cette architecture. Désigne tout outil d'\gls{IA} reposant sur une architecture en couche successive de réseaux neuronaux artificiels, conçu pour imiter le cerveau humain. Cette architecture est ensuite entraînée en autonomie, ou presque, sur de grandes quantités de données. De cette manière, ces outils se spécialisent et développent la capacité de traiter de nouvelles données.}
}

\newglossaryentry{VLM}{
	name={\textit{Vision Language Model}},
	first={\textit{Vision Language Model}},
	text={VLM},
	sort={vision language model},
	description={Modèles d'intelligence artificielle avec des capacités multimodales. C'est à dire que ce modèle n'est plus limité à un champ d'action : il est capable de traiter l'image comme le texte en autonomie.}
}

\newglossaryentry{intero}{
	name={interopérabilité des données}, 
	description={Désigne la capacité à rendre les données d'un système compréhensibles, utilisables et échangeables avec un autre système sans intervention humaine. Cela repose sur des normes et standards permettant de structurer les données de manière uniforme.}
}

\newglossaryentry{SRI}{
	name={\textit{Information Retrieval System}},
	first={\textit{Information Retrieval System}},
	text={SRI},
	sort={information retrieval system},
	description={Logiciels spécialisés sur la recherche d'informations dans des données non structurées. Le plus souvent, ces logiciels sont connectés à une base de données afin de stocker les données dans un index permettant de faciliter l'interrogation de volumes de données massifs. Ils fonctionnent comme un intermédiaire entre l'utilisateur et la base de données : il se charge de l'analyse des données, de les découpées en segments qu'il répartit ensuite dans des n\oe{uds} qu'il indexe.}
}

\newglossaryentry{col}{
	name={collections},
	description={Propre aux bases de données NoSQL orientées documents, comme MongoDB, les collections sont un équivalent des tables SQL. On y stocke des documents, c'est-à-dire des données, de même nature sans schéma particuliers ce qui permet donc une grande souplesse et élasticité.}
}

\newglossaryentry{chps}{
	name={champs},
	description={Dans une base de données NoSQL, le champ est l'unité de base du document. Un champ est une clé, un intitulé, qui contient un ensemble de valeur concernant le document et qui ont toutes le même rôle, la même valeur sémantique. Un champ "dates" contiendra, par exemple, toutes les dates. Chaque document possède plusieurs champs, en fonction des différentes données qu'il stocke. La particularité du champ dans une base de données NoSQL est qu'il est rattaché au document et non à la collection, ainsi plusieurs documents d'une même collection peuvent ne pas avoir exactement les mêmes champs.}
}

\newglossaryentry{ac}{
	name={agent conversationnel},
	plural={agents conversationnels},
	description={L'agent conversationnel n'est pas une technologie mais un outil. Ce mot fait référence à la forme que peut prendre un programme avec lequel interagir sur le principe de l'envoi d'une requête et d'une génération de réponse. En somme, on désigne l'interface par laquelle on interagit avec le programme, qui peut être un \gls{chatbot} ou un \gls{LLM}, que le programme en lui-même. Il y a cependant aujourd'hui une confusion claire dans le langage courant entre \gls{ac} et \gls{LLM}.}
}

\newglossaryentry{ux}{
	name={\textit{User Experience Design}},
	sort={user experience design},
	description={Domaine dans la conception de logiciel informatique visant à rendre l'usage de l'application intuitif, facile et accessible pour tout humain. Ce domaine définit un ensemble de règles et de méthodes à utiliser pour remettre l'humain au centre du logiciel}
}

\newglossaryentry{ui}{
	name={\textit{User Interface Design}},
	sort={user interface design},
	description={Domaine dans la conception de logiciel informatique visant à rendre l'interface générale de l'application simple et attractive afin d'en simplifier l'usage applicatif par un humain. Cela s'accompagne d'un ensemble de règles et de méthodes évolutives afin de créer des interfaces utilisables.}
}

\newglossaryentry{POS}{
	name={\textit{Part of Speech}},
	first={\textit{Part of Speech}},
	text={POS},
	sort={part of speech},
	description={Technologie d'\gls{IA} composant une partie du sous-domaine des technologies d'intelligence artificielle : le \gls{NLP}. Le \textit{Part of Speech}, ou POS, c'est l'analyse grammaticale d'une phrase qui aboutit à l'étiquetage du rôle de chaque mot dans une phrase. Chaque mot est analysé au sein de la phrase, son contexte, et est catégorisé dans l'une des classes grammaticales. La décomposition grammaticale permet ainsi d'en faire sortir le sens}
}

\newglossaryentry{NER}{
	name={\textit{Name Entity Recognition}},
	first={\textit{Name Entity Recognition}},
	text={NER},
	sort={name entity recognition},
	description={Technologie d'\gls{IA} composant une partie du sous-domaine des technologies d'intelligence artificielle : le \gls{NLP}. Le \textit{Name Entity Recognition}, ou NER, est la capacité de reconnaître une chaîne de caractère signifiante comme un nom de personne, de lieux, de marque, d'événements ou encore les dates.}
}

\newglossaryentry{Vecsion}{
	name={vectorisation},
	description={Méthode informatique de calcul et d'attribution d'un identifiant informatique unique à une donnée en fonction de certains facteurs clés. C'est-à-dire qu'en amont du processus nous allons lui donner une donnée, comme un mot, et en sortie nous allons obtenir un identifiant unique sous forme de suite décimale. De cette manière, nous évitons toute confusion entre des données similaires mais non identiques.}
}

\newglossaryentry{vcs}{
	name={vecteur},
	description={Issus d'une vectorisation, le vecteur est une suite décimale unique calculée à partir des caractéristiques uniques d'une donnée. Ils peuvent être plus ou moins complexes et donc éviter des confusions. Ces vecteurs calculés sont ensuite stockés dans une base de données vectorielles permettant la comparaison des vecteurs.}
}

\newglossaryentry{LLM}{
	name={\textit{Large Language Model}},
	sort={large language model},
	first={\textit{Large Language Model}},
	text={LLM},
	description={Souvent abrégé par LLM, les \textit{Large Language Model}, ou grand modèles de langues en français, sont des \glspl{IA} basés sur des \gls{tfmrs}. Ils sont spécialisés dans la compréhension de requête, l'analyse de données textuelles, le raisonnement et la génération de réponse. Pour cela, elles sont entraînées sur de très grandes quantités de données par \gls{DPL} afin de leur conférer ces capacités. Ces \gls{IA} sont concentrées autour de l'alliance entre les technologies regroupées et développées au sein du \gls{NLP} afin de pouvoir effectuer de tel traitement.}
}

\newglossaryentry{RAG}{
	name={\textit{Retrieval Augmented Generation}},
	first={\textit{Retrieval Augmented Generation}},
	text={RAG},
	sort={retrieval augmented generation},
	description={Ici nous parlons d'une manière d'utiliser l'\gls{IA} parce que l'on appelle un \textit{pipeline} ici appelé \gls{RAG}. Il consiste à consolider la base de connaissance d'un \gls{LLM} en le connectant à une source de données contrôlées sur laquelle se baser pour la génération de réponses suite à une requête. Il fonctionne en 3 étapes : l'étape de \textit{retrieval} qui consiste en la recherche des données pertinentes afin de répondre à la question utilisateur. L'étape de l'\textit{augmented} qui consiste à enrichir la réponse générée en faisant référence directement aux sources utilisées et identifiées lors de l'étape précédente. Puis l'étape de \textit{génération} qui consiste en la rédaction d'une réponse prenant pour contexte la requête utilisateur.\footnote{Pour plus d'informations sur ce procédé, consulter \footcite{tuzaInterrogerMemoireDentreprise2025}}}
}

\newglossaryentry{chck}{
	name={\textit{Chunking}},
	sort={chunking},
	description={Procédé informatique fortement utilisé dans les différents domaines de l'\gls{IA}, qui consiste à diviser une information en unités plus petites afin de la stocker plus facilement et d'être plus précis notamment dans des processus de \gls{RAG}.}
}

\newglossaryentry{RR}{
	name={\textit{Re-Ranking}},
	sort={re-ranking},
	description={Procédé informatique s'appuyant sur une \gls{IA} généralement placée à la fin d'un processus de type \gls{RAG}. Il permet de réévaluer et réorganiser les résultats obtenus à l'issus du traitement avant transmission à l'utilisateur afin de s'assurer de leur pertinence.}
}

\newglossaryentry{tk}{
	name={\textit{token}},
	sort={token},
	description={En informatique, le \gls{tk} désigne une unité de stockage destinée à contenir une information de taille variable. En \gls{IA}, on s'en sert notamment afin de contenir des mots ou des morceaux de mots et de phrases. Ils désignent ainsi aussi bien ce qui est envoyé par l'utilisateur lors d'une requête que ce qui est généré par l'\gls{IA}.}
}

\newglossaryentry{pipagreg}{
	name={pipeline d'agrégation},
	plural={pipelines d'agrégation},
	description={Spécifique au langage \gls{NSQL}, notamment le système MongoDB, désigne les étapes successives de traitement effectuées sur les données afin de retourner à l'utilisateur les bonnes données. Il peut s'agir de calculs, de filtres ou bien de modifications afin d'homogénéiser le traitement et le résultat.}
}

\newglossaryentry{LEM}{
	name={lemmatiser},
	description={Procédé appliqué dans le domaine du \gls{NLP}, il désigne le fait ramener un mot conjugué, accordé ou décliner à son état de base que l'on appelle le lemme. Un verbe conjugué sera ainsi ramené à son infinitif par exemple. On l'utilise afin d'épurer le texte et d'en faciliter sa \guillemet{compréhension} sémantique}
}

\newglossaryentry{tfmrs}{
	name={transformers},
	description={Désigne une architecture neuronale développée par Google en 2017 et mise très tôt en application dans le domaine du \gls{NLP}\footcite{devlinBERTPretrainingDeep2019}. Ce modèle permet ce que l'on appelle l'\gls{autat} afin de saisir le sens global d'une phrase, d'un texte, et permet la parallélisation des tâches ce qui accélère le traitement et les tâches d'apprentissage. Il vient remplacer les modèles basés sur \gls{CNN} dans le domaine du \gls{NLP} par son net avantage technologique permettant de fiabiliser les résultats.}
}

\newglossaryentry{2SQL}{
	name={\textit{Text-to-SQL}},
	sort={text-to-sql},
	description={Procédé informatique basé sur l'\gls{IA}, le \gls{RAG} et la recherche intelligente. Il consiste faciliter l'interrogation d'une base de données SQL par l'utilisateur à travers l'usage d'un \gls{LLM}. Il fonctionne en deux temps : d'abord la phase de compréhension de la base de données, permettant de construire le \gls{RAG} puis la phase de requêtage. L'idée générale est d'augmenter l'interrogeabilité de la base de données par la compréhension de la base de données et des requêtes lui étant adressées notamment grâce aux \gls{vcs}. \footnote{Pour plus d'informations, voir \footcite{TexttoSQLInterrogerVos2026}}}
}

\newglossaryentry{2NSQL}{
	name={\textit{Text-to-NoSQL}},
	sort={text-to-nosql},
	description={Adaptation du procédé \gls{2SQL} aux particularités du langage et des bases de données \gls{NSQL}. Encore en développement, l'objectif est de permettre de conférer à un \gls{LLM} la capacité de comprendre et d'interragir avec une base de données \gls{NSQL} et ses particularités comme les \glspl{pipagreg} et la semi-structuration du langage. Il se base lui aussi sur un \gls{RAG}.}
}

\newglossaryentry{AGSQL}{
	name={agent SQL},
	description={Désigne un \gls{LLM}, souvent sous forme d'un \gls{ac}, capable de comprendre, interagir et générer du code SQL et de bases SQL. Dans des processus de \gls{2SQL}, il permet notamment de générer et vérifier la requête SQL générée avant de l'exécuter.}
}

\newglossaryentry{DDL}{
	name={\textit{data definition language}},
	sort={data definition language},
	description={C'est un composant du langage SQL constituant à lui seul un langage à part entière. C'est par ce langage que l'on peut créer une base de données SQL ou en modifier sa structure par différentes opérations appelées CRUD pour \textit{Create}, \textit{Read}, \textit{Update} et \textit{Delete}}
}

\newglossaryentry{NSQL}{
	name={NoSQL},
	description={Acronyme de \textit{Not Only SQL}, désigne un ensemble de bases de données, de typologies variées, ne fonctionnant pas sur un système de table comme le SQL. Ces bases de données ont leurs propres manières de structurer les données, de manière beaucoup plus souple et variable, et ont leurs propres manières d'interroger les données stockées. Les différentes typologies sont : les bases clés-valeurs comme Redis, les bases orientées documents comme MongoDB, les bases orientées colonnes comme Cassandra et les bases orientées graphes comme Neo4j.}
}

\newglossaryentry{chatbot}{
	name={chatbot},
	description={Désigne tout d'abord une réalité technique : un programme informatique conçu sur un système de règles permettant de répondre à un utilisateur avec des réponses pré-définies. Se base sur les premières expérimentations de \gls{NLP}, ne correspondant plus à la réalité technique actuelle. Aujourd'hui, ce terme voit un nouvel usage, plus large, puisque désignant tout programme, \gls{IA} ou non, permettant d'interagir avec l'utilisateur par le texte. \footnote{Pour plus d'informations sur cette transition, consulter : \pageref{chatbot}}}
}

\newglossaryentry{AGIA}{
	name={agent IA},
	plural={agents IA},
	description={Contrairement à l'\gls{ac}, l'\gls{AGIA} est un programme basé sur l'\gls{IA} capable d'interagir avec son environnement informatique. Il est capable de prendre des décisions mais aussi d'utiliser des outils informatiques qui lui sont mis à disposition afin d'exécuter des tâches comme la manipulation de données.}
}

\newglossaryentry{FIAF}{
	name={Fédération Internationale des Archives du Film},
	first={Fédération Internationale des Archives du Film},
	text={FIAF},
	description={Fondée en 1938, fédération à portée internationale qui a pour objectif de protéger et d'assurer la préservation de la culture cinématographique, et plus largement de l'\gls{imgani}. Aujourd'hui, elle assume également un rôle de guide dans la standardisation et l'homogénéisation des pratiques dans la gestion des collections d'\gls{imgani}.}
}

\newglossaryentry{FRBR}{
	name={FRBR},
	description={Défini par l'\gls{IFLA} en 1997, c'est un modèle conceptuel conçu pour améliorer la gestion des collections bibliographiques dans les bibliothèques. Cela s'inscrit dans le mouvement de transition numérique de la gestion des collections et insiste sur la création de liens entre les exemplaires. Il repose essentiellement sur une architecture précise : le \gls{WEMI}. Depuis, ce modèle est adopté par de plus en plus d'institutions notamment celles conservant des \glspl{imgani}.}
}

\newglossaryentry{WEMI}{
	name={WEMI},
	description={Fondement du modèle \gls{FRBR}, il s'agit d'un découpage en 4 entités d'une \oe{uvre} artistique y compris littéraire. On y distingue le \gls{oe}, \textit{Work}, qui désigne le concept. L'\gls{exp} qui caractérise une réalisation de ce concept sous une forme définie comme une pièce de théâtre, ou un film. La \gls{man} qui désigne toute production découlant de l'\gls{exp}. Puis l'\gls{it} qui désigne l'exemplaire que l'on conserve. Cette déconstruction en 4 entités permet de créer plus de liens au sein des catalogues et donc des recherches transversales.}
}

\newglossaryentry{ont}{
	name={ontologie},
	plural={ontologies},
	description={Désigne un moyen informatique d'organisation d'un domaine de connaissance, utilisé dans le Web Sémantique afin de créer des liens. L'ontologie se divise en plusieurs entités : la classe qui regroupe les objets et idées du domaine de connaissance choisi. La propriété qui est une caractérisitique attribuée à une classe. Des relations qui sont caractérisées permettant de faire des liens entre les concepts et d'en expliquer la nature. Et des axiomes qui sont des ensembles de règles pour précisions du sens des termes.}
}

\newglossaryentry{vrt}{
	name={variante},
	description={Notion propre au monde de l'\gls{imgani} qui désigne une version divergente d'une \oe{uvre} originale sans constituer une \oe{uvre} à part entière. Introduit par la \gls{FIAF} dans le \gls{WEMI}, il permet de circonscrire les modifications que peut vivre une \oe{uvre} cinématographique durant son cycle de vie qui viennent en modifier la structure ou le contenu sans que cela soit drastique.}
}

\newglossaryentry{sq}{
	name={séquence},
	description={Dans le milieu de l'\gls{imgani}, désigne un regroupement de \gls{sc} dans une unité de temps basée sur le partage d'une thématique commune, de sujet commun et contribuant à l'avancée d'un même point d'une intrigue. Notion surtout utilisée dans le domaine cinématographique, appartenant à la structure basique d'un film.}
}

\newglossaryentry{sc}{
	name={scène},
	description={Dans le milieu de l'\gls{imgani}, désigne l'unité la plus importante, notamment, d'un film. Une scène est définie comme un ensemble de \gls{pl} faisant figurer les mêmes personnages, dans les mêmes lieux, autour des mêmes actions et sujets. Notion surtout utilisée dans le domaine cinématographique, mais également applicable à toute \gls{imgani}}
}

\newglossaryentry{pl}{
	name={plan},
	description={Dans le milieu de l'\gls{imgani}, désigne une prise de vue simple, qui peut être fixe ou en mouvement, faisant figurer des personnes et/ou des lieux et des actions. Elle se compose de plusieurs images, et peuvent être plus ou moins courtes en fonction des choix de réalisation. C'est l'unité la plus petite d'une \gls{imgani}. Chaque changement d'angle entraînant une nouvelle prise de vue amène un nouveau plan.}
}

\newglossaryentry{table}{
	name={table},
	description={La table constitue la structure de base d'une base de données SQL. Constituée de colonnes et de lignes, elle permet de contenir des ensembles de données structurées. On parle d'enregistrement pour la donnée située sur la ligne, et d'attribue pour les colonnes. \newline Dans l'application SkinSoft, une table correspond à une unité thématique de stockage de données. On parle par exemple de la \guillemet{table des archives}, car c'est dans cette partie de l'application que nous stockons les objets de collections que sont les archives. Ces tables sont rattachées à des index \gls{ES}, c'est-à-dire une structure de données particulières dédiée à l'objet concerné par la table.}
}

\newglossaryentry{grdr}{
	name={grille de recherche},
	description={Notion propre à l'application SkinSoft. La \gls{grdr} constitue un moyen simple d'afficher les données souhaitées des objets recherchés au sein de l'application. Hautement modifiable, l'utilisateur choisit quelle donnée il souhaite visualiser.}
}

\newglossaryentry{rush}{
	name={\textit{rush}},
	sort={rush},
	description={Dans le milieu de l'\gls{imgani}, désigne les \gls{pl} qui sont enregistrés au moment de leurs prises de vues. Les \textit{rush} sont par définition bruts, sans modifications effectuées en post-production.}
}

\newglossaryentry{thesaurus}{
	name={thésaurus},
	description={Les thésaurus sont des hiérarchies de termes représentant des concepts rattachés à des domaines. De cette manière, ils constituent des listes de vocabulaires contrôlés fixant ainsi la terminologie utilisable pour décrire certains domaines. Particulièrement utilisé dans le monde patrimonial et les \gls{sgc} afin de fixer la terminologie utilisée dans l'indexation.}
}

\newglossaryentry{mmlts}{
	name={modèles multimodaux},
	description={Modèles d'\gls{IA} permettant d'associer plusieurs domaines de traitement au sein d'un même outil afin de traiter plusieurs types de données. Ils associent généralement un \gls{LLM} à un encodeur dédié à un autre type de données. Ils permettent notamment d'analyser des vidéos ou des sons et de générer à partir de là, le plus fréquemment, des analyses textuelles.}
}

\newglossaryentry{man}{
	name={manifestation},
	description={Niveaux \gls{WEMI} qui désigne toute forme concrète prise par \gls{exp}. Une manifestation correspond à une matérialité bien définie, que l'on peut dater au moins sur une période de production, obéissant à des critères stricts tels que la durée, le nombre de pages, le lieu de production, la maison d'édition, les réalisateurs, etc.}
}

\newglossaryentry{it}{
	name={item},
	description={Niveaux \gls{WEMI} le plus petit, il désigne l'exemplaire physique que l'on conserve au sein d'une institution ou que l'on détient. Il a des caractéristiques propres à sa vie en tant qu'objet indépendant du reste de son \gls{WEMI}.}
}

\newglossaryentry{exp}{
	name={expression},
	description={Niveaux \gls{WEMI} qui désigne une réalisation concrète d'un concept exprimé au niveau \gls{oe}. Sans avoir encore de forme physique très définie, ici on définit une forme intellectuelle qui incarne l'idée intellectuelle exprimée plus tôt.}
}

\newglossaryentry{oe}{
	name={\oe{uvre}},
	sort={oeuvre},
	description={Niveaux \gls{WEMI} le plus haut, il désigne une idée artistique et intellectuelle que l'on réalise par la suite dans une \gls{exp}, une \gls{man} et un \gls{it} qui est le témoin de sa production. C'est le niveau le plus conceptuel, qui, en dehors du modèle adopté par la \gls{FIAF}, n'a pas de matérialité physique.}
}

\newglossaryentry{hist}{
	name={histogramme},
	description={La définition accordée à ce terme ici est celle de l'histogramme d'une bande vidéo. Désigne le graphique destiné à représenter la distribution de la lumière, de l'ombre et des tons au sein d'une image notamment une \gls{imgani}. Permet d'analyser la composition lumineuse d'une image.}
}

\newglossaryentry{ft}{
	name={\textit{fine-tuning}},
	sort={fine-tuning},
	description={Technique en \gls{IA} qui consiste au réentraînement d'un modèle, comme un \gls{LLM}, sur un nouveau jeu de données. Cela afin de le spécialiser sur une tâche précise, pour laquelle il n'est pas à l'origine. En l'entraînant sur un nouveau jeu de données, le modèle réajuste ces paramètres afin de se spécialiser sur le jeu de données en question. Un \gls{ft} trop poussé peut conduire à une sur-spécialisation du modèle, devenant efficace uniquement sur le jeux de données en question.}
}

\newglossaryentry{CNN}{
	name={réseau de neuronnes convultifs},
	first={réseau de neuronnes convultifs},
	text={CNN},
	description={C'est un type de réseaux de neurones créé pour l'analyse d'images et de vidéos. Chaque couche de réseaux de neurones applique un traitement particulier à l'image ou à la bande vidéo traitée. Chaque couche est dédiée à un traitement, qui s'effectue de manière successive. Nous avons généralement trois couches : la couche dite de \guillemet{convultion} qui déconstruit l'image en zone et en extrait des motifs simples d'abord et complexes ensuite. La couche de regroupement qui sert à compresser les données afin d'accélérer et simplifier le traitement. La couche de connexion qui détermine le résultat final. Ces modèles permettent de faire de la \gls{cl}, de la \gls{seg} et de la \gls{do}.}
}

\newglossaryentry{CLIP}{
	name={CLIP},
	description={Développé par Open-AI, il s'agit d'un \gls{mmlts} capable de comprendre une image dans son sens global et de la décrire avec des mots. Ce modèle est aujourd'hui utilisé dans la plupart des \gls{VLM} afin de faciliter le traitement des images et \gls{imgani}.}
}

\newglossaryentry{ops}{
	name={\textit{open-source}},
	sort={open-source},
	description={Désigne un mouvement informatique de mise à disposition du public le code informatique d'un logiciel. Cela permet donc de récupérer ces logiciels ou \gls{mod} d'\gls{IA}, de les modifier et de les utiliser au sein de \textit{pipelines} ou d'autres logiciels. Ce mouvement s'oppose aux logiciels dits propriétaires, gardant le code logiciel secret et ne permettant pas de le modifier. Cependant, \gls{ops} ne signifie pas gratuité et liberté absolue. Les logiciels appartenants à ce mouvement sont encadrés par des licences, réglementant leurs usages et distributions et peuvent être payants.}
}

\newglossaryentry{autat}{
	name={auto-attention},
	description={Méthode courante du \gls{NLP} permettant à une \gls{IA} d'évaluer l'importance de mots au sein d'une phrase. Cette méthode permet de dépasser le cadre grammatical afin de comprendre les relations conceptuelles entre les mots et le contexte d'une phrase, permettant ainsi de déterminer quel mot est plus déterminant qu'un autre.}
}

\newglossaryentry{GPU}{
	name={processeur graphique},
	description={Composant informatique spécialisé dans le traitement des images et dédié aux calculs à effectuer pour les afficher.}
}

\newglossaryentry{zsht}{
	name={\textit{zero-shot}},
	sort={zero-shot},
	description={En \gls{IA} désigne la capacité d'un \gls{mod} à reconnaître et traiter un objet ou une donnée sur laquelle il n'a pas été entraîné dès la première tentative. Pour cela, le \gls{mod} mobilise toutes ces connaissances afin de déterminer une méthode de traitement adaptée, la mettre en place et la sauvegarder pour des traitements futurs.}
}

\newglossaryentry{diaz}{
	name={diarisation},
	description={Processus informatique basé sur de l'\gls{IA} consistant à pouvoir détecter au sein d'une \gls{imgani} qui parle et quand. De cette manière, on peut associer un dialogue au bon personnage.\footnote{Pour plus d'information, consulter : \url{https://sonix.ai/resources/fr/diarisation-des-locuteurs/}}}
}

\newglossaryentry{CPU}{
	name={processeur},
	description={Composant informatique agissant comme le centre de calcul de l'ordinateur. Il est chargé de réaliser tous les calculs nécessaires afin de lancer les applications, exécuter du code ou encore mobiliser des ressources comme la \gls{RAM} de manière intelligente pour des besoins précis. Il se compose de plusieurs caractéristiques essentielles : les c\oe{urs}, qui sont des composants physiques permettant de paralléliser les tâches, les threads, composant virtuels permettant la même chose, et la fréquence qui détermine la vitesse d'exécution des calculs.}
}

\newglossaryentry{ALM}{
	name={\textit{Audio Large Language Model}},
	first={\textit{Audio Large Language Model}},
	text={ALM},
	sort={audio large language model},
	description={Sous typologie des \gls{mmlts}, elle est spécialisée sur le traitement du texte et des bandes sons. Elle permet d'analyser, transcrire, décrire une bande son.}
}

\newglossaryentry{mmn}{
	name={multimodalité native},
	plural={multimodalités natives},
	description={Dans l'\gls{IA}, désigne un système conçu pour traiter tout type de données en entrée et créer tout type de donnée en sortie. Plus précisément, ce système doit être capable de traiter des données vidéos, sonores et textuelles en même temps, de la même manière et en conservant les liens les unissant si tel est le cas afin d'améliorer sa capacité d'analyse et de compréhension. Ce principe souhaite révolutionner les \gls{mmlts} actuels, ne prenant en charge qu'un seul type de données en entrée et en sortie. Cela pourrait révolutionner la compréhension de documents composites et complexes comme les \glspl{imgani}}
}

\newglossaryentry{NVM}{
	name={\textit{native multimodal modeling}},
	sort={native multimodal modeling},
	description={Désigne les modèles d'\gls{IA} correspondants aux principes de \gls{mmn}. Ils couvrent des modalités et typologies de modèles différentes, obéissant aux mêmes grandes logiques de traitement tout en divergeant un peu les uns des autres. \footnote{Pour plus d'information, voir : \pageref{revNVM} et \footcite{anNativeMultimodalModeling2026}}}
}

\newglossaryentry{sgc}{
	name={système de gestion des collections},
	plural={systèmes de gestion des collections},
	description={Logiciels particulièrement utilisés dans le monde patrimonial permettant de gérer des collections d'objets patrimoniaux de manière informatique. Ces fonctions principales sont : le catalogage, le suivi, la gestion de l'inventaire et la description des collections numériques et physiques. Ce sont des logiciels aujourd'hui très complexes et complets se comportant comme de véritables outils de gestion et de prise de décision stratégique sur la gestion des collections patrimoniales. Ils constituent le \coeur des processus de traitement des collections patrimoniales et sont au centre des pratiques et enjeux métiers. Pour plus d'informations, voir \fullref{SGC?}}
}

\newglossaryentry{IAG}{
	name={intelligence artificielle générative},
	first={intelligence artificielle générative},
	text={IA générative},
	plural={intelligences artificielles génératives},
	description={Typologie de modèle d'\gls{IA}, généralement des \gls{LLM}, capable de générer, à partir d'une requête, du texte, des images, des vidéos ou des sons. Ils ne se contentent pas d'analyser des données, mais de générer quelque chose qu'on leur a demandé à partir de ces analyses. Ces modèles sont aujourd'hui extrêmement répandus et utilisés dans plein de domaines. Les \gls{ac} publics comme ChatGPT, Mistral et Claude sont des \glspl{IAG}}
}

\newglossaryentry{coref}{
	name={coréférence},
	description={Méthode de traitement propre au \gls{NLP} permettant d'associer un mot, une entité, à toutes ses références au sein d'une phrase ou d'un texte entier.}
}

\newglossaryentry{OM}{
	name={\textit{opinion mining}},
	sort={opinion mining},
	description={Méthode de traitement propre au \gls{NLP} avec pour objectif d'identifier, extraire et analyser les sentiments, émotions et prises de positions exprimées au sein d'un texte. Permet l'analyse de texte dans le contenu intellectuel, notamment sollicité dans l'analyse de discours.}
}

\newglossaryentry{seg}{
	name={segmentation},
	description={Dans le domaine de la \gls{CV} désigne la capacité d'une \gls{IA} de diviser l'image en fonction des objets et entités qu'elle y trouve. Sans les identifier précisément, la segmentation permet d'isoler les objets et entités présentes les unes des autres et donc de resituer la composition d'une image.}
}

\newglossaryentry{do}{
	name={détection d'objet},
	description={Dans le domaine de la \gls{CV}, désigne la capacité d'une \gls{IA} de repérer et identifier un objet précis, pour lequel elle a été entraînée, au sein d'une image. On confonds souvent \gls{seg} et \gls{do}. Contrairement à la \gls{seg}, il ne s'agit pas de reconnaître tous les objets et toutes les entités au sein d'une image.}
}

\newglossaryentry{cl}{
	name={classification},
	description={Dans le domaine de la \gls{CV}, désigne la capacité d'une \gls{IA} à analyser une image et de lui attribuer une étiquette agissant comme une identité globale. Elle rattache des images à des groupements déjà prédéfinis selon des caractéristiques prédéfinies également. De cette manière, un \gls{mod} peut faire la distinction entre une photo, une peinture et une gravure par exemple.}
}

\newglossaryentry{mod}{
	name={modèle},
	description={Désigne tout programme informatique basé sur de l'\gls{IA} afin d'apprendre et réaliser des tâches, des prédictions ou bien des décisions à partir de données issues d'une requête. Le \gls{mod} constitue l'outil de base en ingénierie d'\gls{IA}.}
}

\newglossaryentry{RAM}{
	name={mémoire vive},
	description={Composant informatique qui définit la puissance de calcul et la rapidité de leur exécution. Plus elle est importante, plus il sera possible de réaliser un grand nombre de calculs plus rapidement et donc de traiter plus de données.}
}

\newglossaryentry{VRAM}{
	name={mémoire vidéo},
	description={Liée aux \gls{GPU}, c'est-à-dire que ces derniers en déterminent la quantité et la disponibilité, elle définit la puissance de calcul dédiée aux traitements des images, vidéos et objets 3D disponibles. Plus il y a de mémoire vidéo, plus le nombre de calculs réalisables dans un court délai appliqués à des images, vidéos et objets 3D augmente. Elle peut aussi être utilisée pour traiter d'autres catégories données, bien que ce ne soit pas son utilisation première et privilégiée.}
}

\newglossaryentry{2ES}{
	name={\textit{Text-to-ElasticSearch}},
	sort={text to elastisearch},
	description={Processus similaire aux \gls{2SQL} et \gls{2NSQL} qui permet d'associer des outils d'\gls{IA} à \gls{ES} afin d'améliorer le système de recherche de ce dernier. Approche encore expérimentale, elle hybride les processus innovants \gls{2SQL} à la sécurité et robustesse de logiciels largement implémentés que sont les \gls{SRI}}
}

\newglossaryentry{ES}{
	name={ElasticSearch},
	description={ElasticSearch est un logiciel appartenant à la famille des \gls{SRI}. Il permet de stocker, analyser et redistribuer de la donnée textuelle sur de grands volumes de données grâce à son système d'index en document \J. C'est une solution \gls{ops} et hautement scalable avec laquelle il est possible d'interagir par des APIs.}
}

\newglossaryentry{fctcall}{
	name={\textit{function-calling}},
	sort={function-calling},
	description={Désigne la capacité accordée à un processus informatique de réutiliser un même bloc code afin d'effectuer une tâche définie, souvent répétitive. Dans l'\gls{IA}, le \gls{fctcall} amène aussi des notions de reconnaissance contextuelle et de prise de décisions de la part du \gls{mod} : il doit comprendre dans quel cas d'usage il est et appliquer le traitement approprié.}
}

\newglossaryentry{DAM}{
	name={DAM},
	sort={dam},
	description={ Le \textit{Digital Asset Management} est un système informatique permettant de stocker, gérer, modifier et diffuser des fichiers multimédias \textit{(images, vidéos, sons,\dots)}. Son objectif est de centraliser ces médias au même endroit, de proposer des modalités de gestion, création et de modification similaires avant de pouvoir les partager. \newline Dans la suite logicielle SkinSoft, le DAM est un module qui permet de gérer les fichiers multimédias au sein de l'application et des différents modules.}
}

\newglossaryentry{RGPD}{
	name={RGPD},
	sort={rgpd},
	description={Le Réglement Général sur la Protection des Données est un texte européen promulgué en mai 2018 afin de régir le traitement et la collecte des données sensibles et personnelles en Europe. Il concerne toute institution ou entreprise, privée comme publique, qui stocke et traite des données personnelles d'autres personnes physiques ou morales. En plus d'établir des règles et des cas d'usage, ce règlement met également en place des sanctions en cas de non-respect}
}